%% kapa: exam=1 round=22 year=2011 period=2 form=A subject=accounting subject_ko=회계학 seq=2 number=2
%% source: Q-Net 감정평가사 기출문제 (cst00309) artlSeq=5036790 file=extracted/1008/5036790/2011년 제22회 감정평가사 1차시험 기출문제/2교시.hwp
%% conversion: kordoc 4.24.51 -> tools/gen1.py (auto); review=proofread-gpt-astra; flags=none
\begin{question}{2}
재무제표에 관한 설명으로 옳지 않은 것은?
\begin{choices}
  \item 기업이 경영활동을 청산 또는 중단할 의도가 있거나, 경영활동을 계속할 수 없는 상황에 처한 경우를 제외하고는 계속기업을 가정하여 재무제표를 작성해야 한다.
  \item 상이한 성격이나 기능을 가진 항목은 구분하여 표시하되 중요하지 않은 항목은 성격이나 기능이 유사한 항목과 통합하여 표시할 수 있다.
  \item 일반적으로 재무제표는 일관성 있게 1년 단위로 작성하여야 하나, 실무적인 이유로 52주의 보고기간을 적용할 수도 있다.
  \item 재무제표의 기간별 비교가능성을 제고하기 위하여 재무보고를 할 때 전기와 당기를 비교하는 형식으로 보고해야 한다.
  \item 계속기업의 가정이 적절한지의 여부를 평가할 때 경영진은 적어도 보고기간말로부터 향후 6개월 기간에 대하여 이용가능한 모든 정보를 고려해야 한다.
\end{choices}
\end{question}
\answer{5}
